\condAgenticCogWriter\ reframes long-form writing for AI agents as a process-level control problem. A top-level \texttt{Monitor} repeatedly chooses whether to plan, draft, or review from persistent writing state and delegates the selected process to a role-specific subagent. Across three long-form benchmark subsets, the full system is preferred to single-pass generation, fixed writing stages, and adaptive task planning in direct same-prompt comparisons.

The component analyses narrow but do not identify the source of that advantage: explicit goals and adaptive process ordering show only small or unreliable effects, and the observed writing loop usually follows one planning--drafting--reviewing cycle. These findings support process-organized, stateful execution as a promising design while leaving the causal mechanism open, and show how cognitive writing theory can yield executable, falsifiable system hypotheses.
